\documentclass{article}
\usepackage{graphicx} % Required for inserting images



\begin{document}


\section{Výsledky experimentů}

Tato kapitola shrnuje výsledky automatizovaných experimentů pro čtyři navazující
varianty modelu. Cílem není predikovat reálný průběh konkrétní epidemie, ale
porovnat, jak se chování systému mění při postupném přidávání mechanismů:
úmrtnosti, kapacitního omezení, spotřebovávaných zdrojů, dočasné imunity,
vakcinace a strategií alokace.

Všechny uvedené hodnoty vycházejí z BehaviorSpace běhů agregovaných skriptem
\texttt{scripts/aggregate\_results.py}. Každá konfigurace byla spuštěna
desetkrát jako samostatné BehaviorSpace opakování. Uváděné hodnoty jsou průměry
přes těchto deset opakování; směrodatné odchylky jsou doplněny tam, kde jsou
důležité pro interpretaci variability.

\subsection{Experimentální metodika}

Experimenty jsou rozděleny do dvou skupin. První skupinu tvoří baseline scénáře,
které používají výchozí nastavení jednotlivých modelů a umožňují přímé
srovnání variant \texttt{virus1} až \texttt{virus4}. Druhou skupinu tvoří
rozšířené experimenty, které mění vybrané parametry a sledují citlivost modelu
na konkrétní mechanismy.

\begin{itemize}
\item \textbf{Simulace 1} testuje základní SEIR dynamiku bez úmrtnosti.
\item \textbf{Simulace 2} přidává úmrtnost a kapacitu systému.
\item \textbf{Simulace 3} přidává omezené zdroje, dočasnou imunitu a reinfekce.
\item \textbf{Simulace 4} přidává vakcinaci, rizikové skupiny a strategie alokace.
\end{itemize}

V krátkodobých modelech (\texttt{virus1}, \texttt{virus2}) je hlavní metrikou
doba do stabilizace, tedy počet ticků do okamžiku, kdy v populaci nezůstávají
infekční ani exponovaní agenti. V dlouhodobých modelech (\texttt{virus3},
\texttt{virus4}) běh končí až na limitu 8760 ticků. U těchto modelů proto
finální počet infekčních agentů neznamená jednorázový epidemický vrchol, ale
indikátor perzistence nákazy na konci sledovaného období.

\subsection{Souhrn baseline scénářů}

Tabulka~\ref{tab:baseline_summary} ukazuje základní srovnání všech čtyř
modelů. První dva modely se vždy stabilizují. Třetí a čtvrtý model naopak
zůstávají aktivní až do časového limitu, protože do modelu vstupuje reinfekce
a dlouhodobá zdrojová dynamika.

\begin{table}[htbp]
\centering
\small
\begin{tabular}{l r r r r r}
\hline
\textbf{Model} & \textbf{S} & \textbf{E} & \textbf{I} & \textbf{Úmrtí} & \textbf{Ticků} \\
\hline
Simulace 1 & 0.0 & 0.0 & 0.0 & 0.0 & 1384.6 \\
Simulace 2 & 0.0 & 0.0 & 0.0 & 14.8 & 1353.6 \\
Simulace 3 & 0.0 & 46.2 & 122.4 & 106.4 & 8760.0 \\
Simulace 4 & 0.0 & 35.9 & 105.2 & 37.8 & 8760.0 \\
\hline
\end{tabular}
\caption{Baseline výsledky. Sloupce S, E a I označují finální počty
susceptible, exposed a infected agentů.}
\label{tab:baseline_summary}
\end{table}

Výsledky ukazují tři hlavní závěry. Zaprvé, samotný základní SEIR model vede
při zvolených parametrech k úplnému promoření populace bez zbytkové infekce.
Zadruhé, zavedení úmrtnosti a kapacity systému nemění zásadně dobu do
stabilizace, ale přidává měřitelnou ztrátu populace. Zatřetí, zavedení
reinfekcí a zdrojů mění charakter systému: místo jedné uzavřené vlny vzniká
dlouhodobý stav s přetrvávající infekcí.

\subsection{Simulace 1: základní SEIR}

Baseline scénář v Simulaci 1 skončil úplným promořením populace: všech 500
agentů se v průběhu běhu nakazilo a následně přešlo do stavu recovered. Úmrtí
nejsou v tomto modelu implementována, proto zůstávají nulová. Průměrná doba do
stabilizace byla 1384.6 ticku se směrodatnou odchylkou 40.6 ticku.

Citlivostní experiment měnil pravděpodobnost přenosu
\texttt{transmission-prob}. Výsledky jsou v Tabulce~\ref{tab:s1_transmission}.
Vyšší pravděpodobnost přenosu zkracuje dobu do stabilizace, protože se infekce
rychleji rozšíří celou populací a dříve vyčerpá množinu vnímavých agentů.

\begin{table}[htbp]
\centering
\small
\begin{tabular}{r r r r}
\hline
\textbf{\texttt{transmission-prob}} & \textbf{Ticků} & \textbf{Std. odch.} & \textbf{Finální I} \\
\hline
30 & 1622.9 & 89.7 & 0.0 \\
60 & 1372.8 & 57.5 & 0.0 \\
85 & 1267.5 & 35.3 & 0.0 \\
\hline
\end{tabular}
\caption{Vliv pravděpodobnosti přenosu v Simulaci 1.}
\label{tab:s1_transmission}
\end{table}

Interpretace je důležitá: vyšší přenos zde nevede k vyššímu konečnému počtu
nakažených, protože se při všech třech konfiguracích nakazí celá populace.
Rozdíl je v rychlosti průběhu, nikoli ve finálním rozsahu epidemie.

\subsection{Simulace 2: úmrtnost a kapacita systému}

Simulace 2 zachovává základní epidemickou dynamiku z prvního modelu, ale přidává
úmrtnost a kapacitní omezení. Baseline scénář opět končí bez aktivní infekce,
ale průměrně vzniká 14.8 úmrtí se směrodatnou odchylkou 4.7. Průměrná doba do
stabilizace je 1353.6 ticku, tedy velmi blízko Simulaci 1.

Rozhodující rozdíl je v citlivosti na kombinaci kapacity systému a multiplikátoru
přetížení. Tabulka~\ref{tab:s2_capacity} ukazuje, že nejhorší výsledek vzniká
při nízké kapacitě a vysokém multiplikátoru přetížení. Naopak vysoká kapacita
tlumí dopad i při nepříznivém multiplikátoru.

\begin{table}[htbp]
\centering
\small
\begin{tabular}{r r r r}
\hline
\textbf{Kapacita} & \textbf{Přetížení} & \textbf{Úmrtí} & \textbf{Std. odch.} \\
\hline
80 & 1.2 & 8.7 & 2.3 \\
80 & 3.0 & 23.2 & 3.6 \\
240 & 1.2 & 10.0 & 2.0 \\
240 & 3.0 & 17.7 & 3.6 \\
\hline
\end{tabular}
\caption{Vliv kapacity systému a multiplikátoru přetížení v Simulaci 2.}
\label{tab:s2_capacity}
\end{table}

Při nízké kapacitě zvýšení multiplikátoru z 1.2 na 3.0 zvýší průměrný počet
úmrtí z 8.7 na 23.2. Při vysoké kapacitě je nárůst mírnější, z 10.0 na 17.7.
To podporuje očekávání, že kapacita systému nepůsobí pouze jako samostatný
parametr, ale mění dopad ostatních rizikových faktorů.

\subsection{Simulace 3: zdroje a reinfekce}

Simulace 3 mění chování modelu nejvýrazněji. V baseline scénáři se nakazí celá
populace, ale systém se do konce běhu nestabilizuje. Po 8760 ticích zůstává
v průměru 46.2 exponovaných a 122.4 infekčních agentů. Průměrný počet úmrtí je
106.4, tedy výrazně vyšší než v Simulaci 2.

Příčinou je kombinace dočasné imunity, reinfekcí a zdrojové zátěže. Model již
nepopisuje jednu epidemickou vlnu, ale dlouhodobý proces, ve kterém se část
populace vrací do rizika nákazy a dostupné zdroje ovlivňují závažnost průběhu.

\begin{table}[htbp]
\centering
\small
\begin{tabular}{r r r r r}
\hline
\textbf{Zásoba} & \textbf{Obnova} & \textbf{Cena léčby} & \textbf{Úmrtí} & \textbf{Finální zdroje} \\
\hline
600 & 2 & 6 & 114.7 & 44.4 \\
600 & 2 & 24 & 141.5 & 19.2 \\
600 & 20 & 6 & 107.4 & 598.2 \\
600 & 20 & 24 & 110.3 & 592.8 \\
2600 & 2 & 6 & 113.8 & 26.2 \\
2600 & 2 & 24 & 141.3 & 11.2 \\
2600 & 20 & 6 & 109.1 & 2597.0 \\
2600 & 20 & 24 & 107.0 & 2582.0 \\
\hline
\end{tabular}
\caption{Vliv zásoby, obnovy a ceny léčby v Simulaci 3.}
\label{tab:s3_resources}
\end{table}

Tabulka~\ref{tab:s3_resources} ukazuje, že nejrizikovější nejsou pouze nízké
počáteční zdroje, ale zejména kombinace pomalé obnovy a vysoké ceny léčby.
Konfigurace se zásobou 600, obnovou 2 a cenou léčby 24 vede k 141.5 úmrtím.
Podobně nepříznivě dopadá i vysoká počáteční zásoba 2600, pokud je obnova pomalá
a léčba drahá. To naznačuje, že samotná počáteční zásoba nestačí; pro dlouhý
běh je klíčová průběžná regenerace zdrojů.

\subsection{Simulace 4: strategie, rizikové skupiny a vakcinace}

Simulace 4 přidává heterogenitu populace, vakcinaci a strategie alokace zdrojů.
Baseline scénář stále nedosahuje stabilizace do limitu 8760 ticků, ale proti
Simulaci 3 výrazně snižuje úmrtnost: z 106.4 na 37.8 úmrtí. Finální počet
infekčních agentů zůstává také nižší, 105.2 oproti 122.4.

Výsledky strategického gridu jsou shrnuty v Tabulce~\ref{tab:s4_strategy}.
Nejsilnější efekt má vakcinace. Při nízké vakcinaci 20~\% se počet úmrtí
pohybuje přibližně mezi 74 a 104 podle strategie. Při vysoké vakcinaci 80~\%
klesá typicky pod 22 úmrtí.

\begin{table}[htbp]
\centering
\small
\begin{tabular}{r r r r r}
\hline
\textbf{Strategie} & \textbf{Riziková skupina} & \textbf{Vakcinace} & \textbf{Úmrtí} & \textbf{Finální I} \\
\hline
0 & 10 & 20 & 76.9 & 123.7 \\
0 & 10 & 80 & 13.0 & 80.2 \\
0 & 35 & 20 & 74.6 & 125.0 \\
0 & 35 & 80 & 15.4 & 78.8 \\
1 & 10 & 20 & 99.0 & 107.9 \\
1 & 10 & 80 & 20.8 & 85.9 \\
1 & 35 & 20 & 94.1 & 110.3 \\
1 & 35 & 80 & 19.9 & 71.1 \\
2 & 10 & 20 & 77.8 & 118.2 \\
2 & 10 & 80 & 16.7 & 81.6 \\
2 & 35 & 20 & 74.5 & 124.1 \\
2 & 35 & 80 & 15.6 & 78.6 \\
\hline
\end{tabular}
\caption{Porovnání strategií, podílu rizikové skupiny a vakcinace v Simulaci 4.}
\label{tab:s4_strategy}
\end{table}

Strategie 1 má v těchto experimentech nejvyšší úmrtnost, především při nízké
vakcinaci. Například při 10~\% rizikové skupiny a 20~\% vakcinaci vede k 99.0
úmrtím, zatímco strategie 0 ve stejné konfiguraci vede k 76.9 úmrtím a strategie
2 k 77.8 úmrtím. Strategie 0 a 2 jsou podle této metriky robustnější. Při vysoké
vakcinaci se rozdíly mezi strategiemi zmenšují, protože vakcinace dominuje nad
efektem alokačního pravidla.

\subsection{Porovnání dopadu jednotlivých rozšíření}

Postupné rozšiřování modelu ukazuje, že jednotlivé mechanismy nemění pouze
absolutní hodnoty metrik, ale i povahu sledovaného systému.

\begin{itemize}
\item Přechod ze Simulace 1 na Simulaci 2 přidává úmrtnost, ale zachovává
jednovlnný průběh končící stabilizací.
\item Přechod ze Simulace 2 na Simulaci 3 zavádí dlouhodobou dynamiku, ve které
nákaza přetrvává až do konce běhu.
\item Přechod ze Simulace 3 na Simulaci 4 snižuje úmrtnost i finální počet
infekčních, protože do systému vstupuje vakcinace a řízená alokace zdrojů.
\item Výsledky strategie ukazují, že vysoká vakcinace je v tomto nastavení
silnější stabilizační faktor než samotná volba alokační strategie.
\end{itemize}

Z praktického pohledu je nejdůležitější rozdíl mezi jednorázovou stabilizací
a dlouhodobou perzistencí. Pokud model neobsahuje reinfekce, stačí sledovat
rychlost a finální rozsah epidemie. Jakmile jsou reinfekce povoleny, je nutné
sledovat i stav systému v dlouhém horizontu, protože nulový počet vnímavých
agentů neznamená konec epidemické dynamiky.

\subsection{Vyhodnocení hypotéz}

\textbf{H1: Vyšší intenzita přenosu zrychluje průběh epidemie.}
Tato hypotéza je potvrzena pro Simulaci 1. Při zvýšení
\texttt{transmission-prob} z 30 na 85 klesla průměrná doba do stabilizace
z 1622.9 na 1267.5 ticku. V tomto nastavení se ale nezměnil finální rozsah
nákazy, protože se ve všech testovaných variantách nakazila celá populace.

\textbf{H2: Omezená kapacita a zdroje zvyšují úmrtnost při nepříznivých
konfiguracích.} Hypotéza je potvrzena. V Simulaci 2 způsobila kombinace nízké
kapacity a vysokého multiplikátoru přetížení 23.2 úmrtí, zatímco vysoká kapacita
a nízké přetížení vedly k 10.0 úmrtím. V Simulaci 3 se ukázalo, že pomalá obnova
zdrojů a vysoká cena léčby mohou zvýšit úmrtnost až na 141.5 úmrtí.

\textbf{H3: Strategie a vakcinace mohou stabilizovat systém a snížit ztráty.}
Hypotéza je částečně potvrzena. Simulace 4 proti Simulaci 3 snižuje úmrtnost
z 106.4 na 37.8 v baseline scénáři. V dlouhodobém horizontu ale infekce stále
nezaniká a běh končí až časovým limitem. Vakcinace výrazně snižuje úmrtnost,
zatímco rozdíly mezi strategiemi jsou menší a závisí na konkrétní konfiguraci.

\subsection{Omezení výsledků}

Výsledky je nutné číst jako výstup kontrolované replikační studie, nikoli jako
kalibrovanou epidemiologickou predikci. Model používá zjednodušené kontakty,
syntetickou populaci a parametry zvolené pro porovnatelnost experimentů. U
dlouhodobých modelů navíc dostupné agregace neobsahují skutečný peak infekčních
v čase, ale pouze finální hodnoty po 8760 ticích. Pro detailnější analýzu vln by
bylo vhodné exportovat časové řady a dopočítat maximální zatížení systému,
počet vln a délku období s vyčerpanými zdroji.

Dalším omezením je počet opakování. Deset opakování na konfiguraci stačí pro
základní srovnání trendů, ale pro přesnější kvantifikaci nejistoty by bylo
vhodné počet opakování zvýšit a doplnit intervaly spolehlivosti. V současné
verzi experimentů navíc nejsou explicitně fixované random seedy, takže přesné
číselné hodnoty se mohou mezi opakovanými spuštěními mírně lišit.

\subsection{Shrnutí}

Experimenty ukazují, že postupné přidávání mechanismů výrazně mění výsledné
chování agent-based modelu. Základní SEIR varianta vede k uzavřené epidemické
vlně bez úmrtí. Kapacitní omezení přidává úmrtnost, ale systém se stále
stabilizuje. Zdroje a reinfekce posouvají model do režimu dlouhodobé perzistence
nákazy, ve kterém je klíčová nejen počáteční zásoba, ale hlavně schopnost zdroje
průběžně obnovovat. Strategie a vakcinace následně dokážou výrazně snížit
úmrtnost, přičemž nejsilnějším faktorem je v provedených experimentech vysoká
vakcinace.

Z hlediska cíle práce je podstatné, že modely tvoří srovnatelnou řadu. Každé
rozšíření přidává nový mechanismus, jehož dopad lze pozorovat v metrikách.
Výsledky tím podporují použití agent-based simulace jako vysvětlujícího nástroje
pro analýzu omezených zdrojů a jednoduchých intervenčních strategií.

\subsection{Replikovatelnost}

Postup spuštění a agregace je reprodukovatelný přes připravené skripty pro
BehaviorSpace.
Baseline scénáře lze spustit skriptem \texttt{scripts/run\_all\_baselines.ps1}
nebo \texttt{scripts/run\_all\_baselines.sh}. Rozšířené experimenty lze spustit
skriptem \texttt{scripts/run\_extended\_experiments.ps1} nebo
\texttt{scripts/run\_extended\_experiments.sh}. Agregace výsledků probíhá
skriptem \texttt{scripts/aggregate\_results.py}, který vytváří soubory
\texttt{out/*\_summary.csv} použité v této kapitole.


\end{document}
